Commands for simulating panel data and life-cycle profiles from a Finite Horizon Value Fn problem are given by \\
\textcolor{RoyalBlue}{SimPanelValues=SimPanelValues\_{}FHorz\_{}Case1(InitialDist, Policy, FnsToEvaluate,} \\
\indent \indent \textcolor{RoyalBlue}{FnsToEvaluateParamNames, Parameters, n\_{}d, n\_{}a, n\_{}z, N\_{}j, d\_{}grid, a\_{}grid, z\_{}grid,}\\
\indent \indent \textcolor{RoyalBlue}{pi\_{}z, [simoptions]);} \\
and \\
\textcolor{RoyalBlue}{SimLifeCycleProfiles=SimLifeCycleProfiles\_{}FHorz\_{}Case1(InitialDist, Policy, FnsToEvaluate,} \\
\indent \indent \textcolor{RoyalBlue}{FnsToEvaluateParamNames, Parameters, n\_{}d, n\_{}a, n\_{}z, N\_{}j, d\_{}grid, a\_{}grid, z\_{}grid,}\\
\indent \indent \textcolor{RoyalBlue}{pi\_{}z, [simoptions]);} \\
This section describes the simulations involved, the inputs and outputs, and provides some further info on using these commands.
They are covered together here as their inputs and functionality have much in common.

Both of these commands are intended to be run based on a $Policy$ that has been output by a $ValueFnIter\_{}FHorz\_{}Case1()$ command.
$PType$ versions of these commands also exist that allow for different agent permanent/fixed types.
Many of the inputs are simply the same as those originally used to solve for $Policy$, namely $Parameters$, $n\_{}d$, $n\_{}a$, $n\_{}z$, $N\_{}j$, $d\_{}grid$, $a\_{}grid$, $z\_{}grid$, $pi\_{}z$,
and so these are not redescribed here.

$simoptions$ allows you to set some internal options (including parallization), if $simoptions$ is not used all options will revert to their default values.

The forms that each of these inputs and outputs takes are now described in detail. The best way to understand how to use the
command may however be to just go straight to the examples; in particular the replication of \cite*{HubbardSkinnerZeldes1994}.

\subsection{Inputs and Outputs}
To use the toolkit to solve problems of this sort the following steps must first be made.
\begin{itemize}
 \item To be able to simulate a panel you must first declare in what state agents will be 'born'. This is the
   role of $InitialDist$. Typically it will be a probability distribution across the state space $(a,z)$ at age $j=1$.
   This will take the form of a matrix of dimensions $n\_{}a-by-n\_{}z$ the elements of which sum to one. \\
   \textcolor{RoyalBlue}{InitialDist=zeros(n\_{}a,n\_{}z); InitialDist(1,ceil(n\_{}z)/2)=1;} \\
   (In this example it is assumed agents are born with least number of assets ($n\_{}a$ one-dimensional) and median values of the shocks.) \\
   ($InitialDist$ is allowed to further depend on $j$, in which case it will be $n\_{}a-by-n\_{}z-by-N\_{}j$. When the $N\_{}j$ dimension
   is not declare it is simply assumed by default that all agents are 'born' with $j=1$.)
 \item The policy function $Policy$ is simply that created as an output by $ValueFnIter\_{}Case1\_{}FHorz()$. It contains the 'indexes' for
   the optimal policy (choices for $d$ and $a'$ variables) at each point in the discretized state space.
 \item Define functions for the variables you wish to create panel data (or life-cycle profiles) for.
   Store all these functions together. \\
   eg., for a panel which is just the values of the endogenous state, then \\
   \textcolor{RoyalBlue}{FnsToEvaluateFn\_{}1 = @(aprime\_{}val,a\_{}val,s\_{}val) a\_{}val;} \\
   (You always need the $@(...)$ defined either in this way, or with $@(d\_{}val,...)$ where appropriate. You also need to include parameters where appropriate.) \\
   Define the names of any parameters used to evaluate the aggregate variables. \\
   \textcolor{RoyalBlue}{FnsToEvaluateParamNames(1).Names=\{\};} \\
   (In our example there were none, but an example might be that you want to evaluate tax revenue, which is a tax rate parameter times an endogenous state variable, then the tax rate
    parameter name will be given here.) \\
   You then store all of these together: \\
   \textcolor{RoyalBlue}{FnsToEvaluate=\{FnsToEvaluateFn\_{}1\};}
 \item For how to declare $Parameters$, $n\_{}d$, $n\_{}a$, $n\_{}z$, $N\_{}j$, $d\_{}grid$, $a\_{}grid$, $z\_{}grid$, $pi\_{}z$
   check the descriptions of the finite horizon value function iteration command in Section \ref{sec:VFI_Case1_FHorz}.
\end{itemize}

That covers all of the objects that must be created, the only thing left to do is simply call the commands. \\
\textcolor{RoyalBlue}{SimPanelValues=SimPanelValues\_{}FHorz\_{}Case1(InitialDist, Policy, FnsToEvaluate,} \\
\indent \indent \textcolor{RoyalBlue}{FnsToEvaluateParamNames, Parameters, n\_{}d, n\_{}a, n\_{}z, N\_{}j, d\_{}grid, a\_{}grid, z\_{}grid,}\\
\indent \indent \textcolor{RoyalBlue}{pi\_{}z, [simoptions]);} \\
and \\
\textcolor{RoyalBlue}{SimLifeCycleProfiles=SimLifeCycleProfiles\_{}FHorz\_{}Case1(InitialDist, Policy, FnsToEvaluate,} \\
\indent \indent \textcolor{RoyalBlue}{FnsToEvaluateParamNames, Parameters, n\_{}d, n\_{}a, n\_{}z, N\_{}j, d\_{}grid, a\_{}grid, z\_{}grid,}\\
\indent \indent \textcolor{RoyalBlue}{pi\_{}z, [simoptions]);} \\

The outputs are
\begin{itemize}
 \item $SimPanelValues$: A simulated panel data set containing the values of the 'FnsToEvaluate'. It will be a matrix of size [$length(FnsToEvaluate)$,N\_{}j, $NSims$]
    and at each point it will be the value of the 'FnsToEvaluateFn' for an age and individual,
 \item $SimLifeCycleProfiles$: Creates life-cycle profiles for the values of the 'FnsToEvaluate'. It will be a matrix of size [$length(FnsToEvaluate)$, N\_{}j, 23].
    The first dimension indexes the variable, the second is the age, and the third indexes mean/median/ventile.
    For example, $SimLifeCycleProfiles(3,:,1)$ will be the life-cycle profile for the mean of the third variable in 'FnsToEvaluate'. \\
    \small ($simoptions.lifecyclepercentiles=20$ is ventiles,
    and can be used to control the number of percentiles reported (always including the min and max values), so $simoptions.lifecyclepercentiles=4$
    will return quartiles, and $simoptions.lifecyclepercentiles=0$ then only the mean and median will be returned.)
    \normalsize
\end{itemize}

\subsection{Some further remarks}
\begin{itemize}
 \item Models where $d$ is unnecessary (only $a'$ need be chosen): set $n\_{}d=0$ and $d\_{}grid=0$, the code will
    take care of the rest.
 \item $SimLifeCycleProfiles$ simply performs the same panel data simulation as $SimPanelValues$ but then returns the age-conditional mean (median/percentile/etc.)
    rather than returning the actual panel simulation.
 \item Panel simulation commands presently just ignore stochastic survivial probabilities (the value function and OLG commands do account for them).
\end{itemize}

\subsection{Options}
Optionally you can also input a further argument, a structure called \textit{simoptions}, which allows you to set various internal options.
Following is a list of the \textit{simoptions}, the values to which they are being set in this list are their default values.
 \begin{itemize}
  \item For reasons of speed simulations are always performed on CPU (or parallel CPUs). Setting $simoptions.parallel=2$ simply
    means that the output will be transferred to the GPU before it is returned (and that inputs are assumed to be on
    the GPU). \\
    \textcolor{RoyalBlue}{simoptions.parallel=2}
  \item If you want feedback set to one, else set to zero \\
    \textcolor{RoyalBlue}{simoptions.verbose=0}
  \item Number of periods for which each individual agent simulation lasts. \\
    \textcolor{RoyalBlue}{simoptions.simperiods=N\_{}j}
  \item Number of individual agents simulations that make up the panel data. \\
    \textcolor{RoyalBlue}{simoptions.numbersims=10\^{}4} \\
    (Note: Life-cycle profiles are themselves based on a simulated panel so this $simoptions$ is relevant)
  \item Number of percentiles for which life-cycle profiles are given (not relevant to $SimPanelValues$ command),
    in addition to the mean, median, and minimum (note that maximum will always be the p-th of the p percentiles) \\
    \textcolor{RoyalBlue}{simoptions.lifecyclepercentiles=20;}
\end{itemize}

\subsection{Some Examples}
See the replication of \cite*{HubbardSkinnerZeldes1994} at github: \href{https://github.com/vfitoolkit/vfitoolkit-matlab-replication/tree/master/HubbardSkinnerZeldes1994}{https://github.com/vfitoolkit/vfitoolkit-matlab-replication/tree/master/HubbardSkinnerZeldes1994}.
